\documentclass[11pt,a4paper]{article}

\usepackage[utf8]{inputenc}
\usepackage[T1]{fontenc}
\usepackage{lmodern}
\usepackage{amsmath, amssymb, amsfonts}
\usepackage{bm}
\usepackage{geometry}
\usepackage{graphicx}
\usepackage{hyperref}
\usepackage{authblk}
\usepackage{cite}
\usepackage{physics}
\usepackage{mathtools}

\geometry{margin=2.4cm}

\title{\textbf{Companion XCVII--XCVIII \\ Degeneracy Structure and Null Directions in RDCC}}
\author{Michael Lehmann \\ \texttt{mi.lehmann@gmx.de}}
\date{April 2026}

\begin{document}
\maketitle

\begin{abstract}
This merged Companion unifies the content of Companions XCVII--XCVIII into a 
single document. It presents the degeneracy structure of RDCC inference, the 
geometry of null directions in the Fisher manifold, and the alignment between 
weakly constrained modes across geometric, morphological, and topological 
observables. Together, these components provide a complete description of the 
flat and weakly informative regions of the RDCC likelihood landscape.
\end{abstract}
\section{Degeneracy Structure in RDCC Inference}

Even though RDCC introduces only a single physical parameter, the relaxation 
parameter $\alpha$, the inference landscape exhibits nontrivial degeneracy 
structure. These degeneracies arise from scale-dependent responses of 
observables, partial cancellations between geometric and topological probes, 
and the presence of weakly informative regions in the Fisher manifold.

\subsection{Definition of degeneracy directions}

A degeneracy direction is defined as a direction in parameter space along which 
the likelihood varies only weakly. In the one-parameter RDCC case, this 
corresponds to regions where
\begin{equation}
\partial_\alpha \log \mathcal{L} \approx 0,
\end{equation}
or equivalently where the Fisher information
\begin{equation}
g(\alpha) = \mathbb{E}\!\left[ (\partial_\alpha \log \mathcal{L})^2 \right]
\end{equation}
is small.

These regions correspond to:
\begin{itemize}
    \item weak sensitivity of observables to $\alpha$,
    \item partial cancellations between different probes,
    \item noise-dominated scales,
    \item or intrinsic symmetries of the relaxation kernel.
\end{itemize}

\subsection{Scale-dependent degeneracies}

RDCC exhibits a characteristic pattern of degeneracies:
\begin{itemize}
    \item \textbf{Linear scales} ($k \lesssim 0.05\,h\,\mathrm{Mpc}^{-1}$):  
    weak response to $\alpha$, dominated by cosmic variance.

    \item \textbf{Intermediate scales} ($k \sim 0.1$--$1\,h\,\mathrm{Mpc}^{-1}$):  
    strong breaking of degeneracies due to the relaxation kernel.

    \item \textbf{Deeply nonlinear scales}:  
    partial re-emergence of degeneracies due to noise and baryonic effects.
\end{itemize}

This structure is stable across:
\begin{itemize}
    \item tomographic binning schemes,
    \item combinations of geometric and topological probes,
    \item different noise realizations.
\end{itemize}

\subsection{Cross-probe degeneracies}

Degeneracies may arise not only within a single observable, but also across 
different probes. For example:
\begin{itemize}
    \item geometric probes (power spectra) may show weak sensitivity where
    \item topological probes (persistent homology, transport-based topology)
    remain strongly informative.
\end{itemize}

This cross-probe structure is essential for:
\begin{itemize}
    \item designing optimal combinations of summary statistics,
    \item identifying complementary probes,
    \item constructing robust likelihoods that avoid information loss.
\end{itemize}

The degeneracy structure thus forms the first component of the merged 
XCVII--XCVIII Companion.
\section{Null Directions and Alignment}

Degeneracy directions describe regions of weak sensitivity in the likelihood, 
but a deeper geometric structure emerges when these directions are analyzed in 
terms of \emph{null directions} of the Fisher metric. These correspond to 
directions in observable space along which variations in $\alpha$ induce 
minimal or vanishing response.

\subsection{Definition of null directions}

A null direction is defined by a vector $u$ in observable space such that
\begin{equation}
u^\top \, \partial_\alpha \mathbf{d} \approx 0,
\end{equation}
where $\mathbf{d}$ denotes the vector of summary statistics.

Equivalently, in terms of the Fisher metric,
\begin{equation}
g(\alpha) \approx 0
\quad \Longleftrightarrow \quad
\partial_\alpha \log \mathcal{L} \approx 0.
\end{equation}

Null directions arise when:
\begin{itemize}
    \item the relaxation kernel induces symmetric distortions,
    \item geometric and topological probes respond in opposite ways,
    \item noise dominates the observable combination,
    \item or the observable is intrinsically insensitive to $\alpha$.
\end{itemize}

\subsection{Alignment across probes}

A key feature of RDCC inference is the \emph{alignment} of null directions 
across different classes of observables. Specifically:
\begin{itemize}
    \item geometric probes (power spectra),
    \item morphological probes (peaks, voids, Minkowski functionals),
    \item and topological probes (persistent homology, transport-based topology)
\end{itemize}
share partially aligned null directions.

This alignment implies:
\begin{itemize}
    \item certain combinations of observables are jointly insensitive to $\alpha$,
    \item degeneracies persist even when probes are combined,
    \item but the \emph{orthogonal} combinations remain highly informative.
\end{itemize}

\subsection{Geometric interpretation}

Null directions correspond to flat regions in the Fisher manifold. In natural 
coordinates, these appear as extended plateaus where the likelihood varies only 
weakly. The alignment of null directions across probes indicates that these 
plateaus are not artifacts of a single observable, but reflect intrinsic 
properties of the relaxation dynamics.

\subsection{Implications for inference}

Understanding null directions is essential for:
\begin{itemize}
    \item constructing stable likelihoods,
    \item avoiding overfitting in compressed representations,
    \item designing probe combinations that maximize information,
    \item and identifying observables that break degeneracies most efficiently.
\end{itemize}

The structure of null directions thus forms the second component of the merged 
XCVII--XCVIII Companion.
\section{Cross-Probe Structure of Degeneracies}

The degeneracy and null-direction structure of RDCC inference is not confined 
to a single observable. Instead, it emerges from the interplay between 
geometric, morphological, and topological probes. This cross-probe structure 
reveals how different summary statistics respond to the relaxation parameter 
$\alpha$ and how their sensitivities overlap or diverge.

\subsection{Shared and complementary degeneracies}

Across probes, we observe:
\begin{itemize}
    \item \textbf{Shared degeneracies}:  
    Certain combinations of observables exhibit similar weak responses to 
    $\alpha$, leading to aligned null directions.

    \item \textbf{Complementary degeneracies}:  
    Some probes remain insensitive where others are strongly informative, 
    enabling degeneracy breaking through probe combination.

    \item \textbf{Stable degeneracy patterns}:  
    The structure persists across noise realizations and tomographic schemes.
\end{itemize}

This structure is essential for designing optimal multi-probe inference 
pipelines.

\subsection{Implications for compressed inference}

The cross-probe degeneracy structure informs:
\begin{itemize}
    \item the construction of compressed statistics that avoid null directions,
    \item the identification of probe combinations that maximize information,
    \item the design of likelihoods that remain stable across probes,
    \item and the interpretation of tensions or synergies between observables.
\end{itemize}

\section{Conclusion}

This merged Companion unifies the degeneracy and null-direction structure of 
RDCC inference. Degeneracy directions identify weakly informative regions of 
the likelihood, while null directions reveal deeper geometric structure in 
observable space. Their alignment across probes highlights intrinsic properties 
of the relaxation dynamics and informs the design of robust, multi-probe 
inference strategies.

Together, these components replace Companions XCVII--XCVIII and provide a 
complete description of the flat and weakly constrained regions of the RDCC 
likelihood landscape.

\begin{thebibliography}{9}

\bibitem{RDCC-Flagship}
M.~Lehmann,
\textit{Relaxation-Driven Cyclic Cosmology (RDCC): Flagship Paper},
(2026).

\bibitem{RDCC-Summary}
M.~Lehmann,
\textit{RDCC Summary Paper: Inference Framework and Key Results},
(2026).

\bibitem{RDCC-Observables}
M.~Lehmann,
\textit{RDCC Numerical Fits and Observational Constraints},
(2026).

\bibitem{RDCC-Topology}
M.~Lehmann,
\textit{Topological and Morphological Probes in RDCC},
(2026).

\bibitem{RDCC-Companions-Deg}
M.~Lehmann,
\textit{RDCC Companion Series: Degeneracy Structure and Null Directions},
Companions XCVII--XCVIII (merged), (2026).

\end{thebibliography}

\end{document}
